\bibitem{sa1} Leighton, T. y Rubinfeld, R. (2006). Mathematics for Computer Science: Notes for Recitation 6. MIT. Definiciones de grafos, recorridos y longitud. \url{https://web.mit.edu/neboat/Public/6.042/rec6.pdf}.
\bibitem{sa2} Diestel, R. (2025). Graph Theory, 6.ª ed. Springer, capítulo 10. Definiciones hamiltonianas; referencia conservada de la investigación de SA. \url{https://www.math.uni-hamburg.de/home/diestel/books/graph.theory/preview/Ch10.pdf}.
\bibitem{sa3} Universidad de Waterloo. The Traveling Salesman Problem. Definición del TSP. \url{https://www.math.uwaterloo.ca/tsp/problem/index.html}.
\bibitem{sa4} Kirkpatrick, S., Gelatt, C. D. y Vecchi, M. P. (1983). Optimization by Simulated Annealing. Science 220(4598), 671--680. Fundamento de SA. \url{https://doi.org/10.1126/science.220.4598.671}.
\bibitem{sa5} Hajek, B. (1988). Cooling Schedules for Optimal Annealing. Mathematics of Operations Research 13(2), 311--329. Condiciones de convergencia. \url{https://doi.org/10.1287/moor.13.2.311}.
\bibitem{sa7} Komosinski, M. Optimization: Simulated annealing. Poznan University of Technology. Material docente citado por el documento original para el pseudocódigo y el enfriamiento. \url{https://www.cs.put.poznan.pl/mkomosinski/lectures/optimization/SA-en.pdf}.
\bibitem{sa8} Metropolis, N., Rosenbluth, A. W., Rosenbluth, M. N., Teller, A. H. y Teller, E. (1953). Equation of State Calculations by Fast Computing Machines. Journal of Chemical Physics 21(6), 1087--1092. Antecedente de la aceptación probabilística. \url{https://doi.org/10.1063/1.1699114}.
\bibitem{pso1995} Kennedy, J. y Eberhart, R. (1995). Particle swarm optimization. Proceedings of ICNN'95, vol. 4, 1942--1948. Fuente original de PSO. \url{https://doi.org/10.1109/ICNN.1995.488968}.
\bibitem{pso1998} Shi, Y. y Eberhart, R. (1998). A modified particle swarm optimizer. IEEE International Conference on Evolutionary Computation, 69--73. Variante con peso de inercia. \url{https://doi.org/10.1109/ICEC.1998.699146}.
% --- Añadidas para el reporte de la Fase 4 (verificar los datos antes de entregar) ---
\bibitem{wilson1927} Wilson, E. B. (1927). Probable Inference, the Law of Succession, and Statistical Inference. Journal of the American Statistical Association 22(158), 209--212. Intervalo de confianza usado para la tasa de éxito. \url{https://doi.org/10.1080/01621459.1927.10502953}.
\bibitem{posa1976} Pósa, L. (1976). Hamiltonian circuits in random graphs. Discrete Mathematics 14(4), 359--364. Rotaciones usadas en la variante ACO + Pósa. \url{https://doi.org/10.1016/0012-365X(76)90068-6}.
\bibitem{slh2014} Baniasadi, P., Ejov, V., Filar, J. A., Haythorpe, M. y Rossomakhine, S. (2014). Deterministic ``Snakes and Ladders'' Heuristic for the Hamiltonian cycle problem. Mathematical Programming Computation 6(1), 55--75. Heurística de referencia del grupo FHCP. \url{https://doi.org/10.1007/s12532-013-0059-2}.
\bibitem{lkh2000} Helsgaun, K. (2000). An effective implementation of the Lin--Kernighan traveling salesman heuristic. European Journal of Operational Research 126(1), 106--130. Base de LKH, usada para resolver instancias del HCP como TSP. \url{https://doi.org/10.1016/S0377-2217(99)00284-2}.
